\section{INTRODUCTION}

Fast flat-ground transport and obstacle negotiation impose different mechanical
requirements. Wheels provide continuous contact and efficient rolling, but their
obstacle capability is limited by geometry. Articulated legs can place contacts
and reposition the body, but generally require more actuators, gait coordination,
and time per unit distance. Hybrid platforms therefore face a design question:
how can one robot retain fast transport while crossing stairs without adding a
dedicated wheel drive or a separate transformation mechanism?

This paper investigates that question through \robotname{}, shown schematically
in Fig.~\ref{fig:teaser}. Each of its six three-actuator legs terminates in an
actuated arc-shaped distal link. Adjacent arcs can provide approximately
continuous rolling support on flat ground, while the open arc and articulated
leg joints remain available for walking and stair-edge engagement.

\begin{figure*}[t]
  \centering
  \fbox{\parbox[c][31mm][c]{0.94\textwidth}{\centering
    Final Fig.~1: one clean side view of the robot, followed by three frames of
    battery-powered walking, quasi-rolling, and continuous stair ascent. Include
    scale and anonymize faces, affiliations, logos, and external links.}}
  \caption{The proposed morphology uses the same actuated arc-shaped distal
  links for three locomotion modes. This figure must communicate the complete
  paper claim before the reader reaches the method section.}
  \label{fig:teaser}
\end{figure*}

The central hypothesis is that the shared distal mechanism can reduce the
mobility trade-off between a simple rotating-leg robot and a fully articulated
hexapod. The paper evaluates three claims: (i) the platform can transition among
walking, quasi-rolling, and stair-climbing configurations without mechanical
reassembly; (ii) quasi-rolling improves flat-ground transport relative to
articulated walking on the same hardware; and (iii) coordinated articulation and
arc engagement improve stair traversal relative to a distal-only configuration.

The intended contributions are:
\begin{itemize}
  \item an eighteen-actuator hexapod morphology with six multifunctional,
  actuated arc-shaped distal links;
  \item a mode-transition and locomotion controller that reuses the same hardware
  for walking, quasi-rolling, and stair climbing; and
  \item a controlled experimental comparison using the same platform, including
  speed, success rate, traversal time, body attitude, and electrical demand.
\end{itemize}

\draftnote{After all experiments are frozen, rewrite the final paragraph so that
each contribution ends with one concrete, verified result.}
